\section*{अध्याय \ref{ch:g2:balance-equilibrium} --- संतुलन और साम्यावस्था}

\begin{solution}{exo:g2:balance-equilibrium:1}
पृथ्वी लैंप को नीचे खींचती है; मेज़ उसे ठीक उतने ही ज़ोर से ऊपर
धकेलती है। दोनों बलों में बराबरी हो जाती है --- \omterm{def:g2:balance-equilibrium:equilibrium}{साम्यावस्था} ---
इसलिए लैंप अपनी जगह बना रहता है।
\end{solution}

\begin{solution}{exo:g2:balance-equilibrium:2}
जब दोनों टोलियाँ उलटी दिशाओं में बराबर ज़ोर से खींचती हैं: खिंचाव
एक-दूसरे को काट देते हैं। जिस पल एक टोली दूसरी से ज़्यादा ज़ोर लगाती
है, बराबरी टूट जाती है।
\end{solution}

\begin{solution}{exo:g2:balance-equilibrium:3}
वह सीधा टिका रहता है --- बराबर दूरियों पर बराबर भार। जब एक जुड़वाँ और
दूर खिसक जाता है, तब उसी की ओर नीचे जाती है: दूरी भी गिनती में आती
है।
\end{solution}

\begin{solution}{exo:g2:balance-equilibrium:4}
धुरी के पास, और लीना अपनी ओर दूर सिरे पर। नियम: बीच के पास बैठा भारी
सवार दूर बैठे हल्के सवार से बराबरी कर सकता है --- भार और दूरी, दोनों
गिनती में आते हैं।
\end{solution}

\begin{solution}{exo:g2:balance-equilibrium:5}
रूलर का \omterm{def:g2:balance-equilibrium:point}{संतुलन बिंदु} उसके बीच में है; हथौड़े का भारी सिर के पास।
जाँच: हर एक को एक उँगली पर टिकाकर वह जगह ढूँढ़ो जहाँ वह सीधा रहता
है।
\end{solution}

\begin{solution}{exo:g2:balance-equilibrium:6}
वे बुहारी के पास मिलती हैं --- यानी भारी सिरे की ओर। किसी वस्तु का
\omterm{def:g2:balance-equilibrium:point}{संतुलन बिंदु} उसके भारी पक्ष की ओर बैठता है।
\end{solution}

\begin{solution}{exo:g2:balance-equilibrium:7}
चौड़ी पिरामिड: वह काफ़ी दूर तक झुक सकती है, तब जाकर उसका भार अपने
चौड़े आधार से बाहर लटकता है। पतली मीनार ज़रा-से झुकाव पर ही गिर जाती
है।
\end{solution}

\begin{solution}{exo:g2:balance-equilibrium:8}
बाँस हर डगमगाहट को धीमा कर देता है, जिससे नट को वापस रस्सी के ऊपर आने
का समय मिल जाता है। फ़ुटपाथ के किनारे पर तुम इसी वजह से बाँहें फैला
लेते हो: डगमगाहट धीमी, सँभलना आसान।
\end{solution}

\begin{solution}{exo:g2:balance-equilibrium:9}
अकेले दाएँ पैर पर खड़े होने के लिए तुम्हारे शरीर को दाईं ओर झुकना
पड़ेगा, ताकि तुम्हारा भार उसी पैर के ऊपर आ जाए --- ठीक वैसे जैसे
मीनार अपना भार अपने आधार के ऊपर रखती है। दीवार ठीक उसी ओर है जिस ओर
झुकना है: वह इस झुकाव को मना कर देती है, तुम्हारा भार उठे हुए पैर की
ख़ाली जगह के ऊपर लटका रह जाता है, और तुम लुढ़क जाते हो।
\end{solution}

\begin{solution}{exo:g2:balance-equilibrium:10}
क्रेन एक विशाल सी-सॉ है और उसका खंभा धुरी। पहला सवार: बोझ, लंबी अगली
भुजा पर दूर। दूसरा सवार: सीमेंट के गुटके --- कहीं ज़्यादा भारी,
इसलिए वे पास बैठ सकते हैं, यानी छोटी पिछली भुजा पर। भारी-पास दूर के
हल्के से बराबरी कर लेता है, ठीक सी-सॉ का नियम; पिछली भुजा इसलिए छोटी
है कि उस पर बैठा सवार भारी है।
\end{solution}
